\documentclass{article}

\usepackage{../../common}

\graphicspath{
  {../../../media/images/}
}
\addbibresource{../../references.bib}

\title{Deliteľnosť prirodzených čísel}
\begin{document}

\maketitle

\section{Deliteľnosť čísel dvojkou}

Prirodzené čísla delíme podľa deliteľnosti dvoma na párne a nepárne. Párne čísla sú také prirodzené čísla, ktoré sú deliteľné číslom \(2\) bez zvyšku. Párne čísla možno rozdeliť na dve rovnaké skupiny. Nepárne čísla sú také prirodzené čísla, ktoré nie sú deliteľné číslom \(2\) bez zvyšku. Pri delení dvoma zostáva vždy zvyšok \(1\).

Párnosť alebo nepárnosť prirodzeného čísla určujeme podľa poslednej číslice (rád jednotiek). Posledná číslica párneho čísla je vždy \(0\), \(2\), \(4\), \(6\) alebo \(8\). Posledná číslica nepárneho čísla je vždy \(1\), \(3\), \(5\), \(7\) alebo \(9\).
\begin{itemize}
  \item Príkladmi párnych čísel sú čísla \(0\), \(2\), \(4\), \(6\), \(8\), \(10\), \(12\), \(14\), \(16\), \(18\), \(20\)\dots
  \item Príkladmi nepárnych čísel sú čísla \(1\), \(3\), \(5\), \(7\), \(9\), \(11\), \(13\), \(15\), \(17\), \(19\), \(21\)\dots
\end{itemize}

Párne číslo môžeme zapísať všeobecne ako \(2n\), kde \(n\) je ľubovoľné prirodzené čislo. Ak za \(n\) dosadíme prirodzené čisla \(0\), \(1\), \(2\), \(3\), \(4\), \(5\), \(6\)\dots dostaneme postupnosť \(0\), \(2\), \(4\), \(6\), \(8\), \(10\), \(12\)\dots Nepárne číslo môžeme zapísať všeobecne ako \(2n+1\), kde \(n\) je ľubovoľné prirodzené čislo. Ak za \(n\) dosadíme prirodzené čisla \(0\), \(1\), \(2\), \(3\), \(4\), \(5\), \(6\)\dots dostaneme postupnosť \(1\), \(3\), \(5\), \(7\), \(9\), \(11\), \(13\)\dots

Párnosť alebo nepárnosť prirodzeného čísla závisí iba od číslice v ráde jednotiek (poslednej číslice). Všetky vyššie rády (desiatky, stovky, tisíce\dots) sú násobkami čísla \(10\), a teda aj násobkami čísla \(2\). Na hodnotu číslice rádu jednotiek súčtu (resp. súčinu) majú vplyv iba hodnoty číslic rádu jednotiek dvoch sčítancov (resp. činiteľov). Pri sčítaní (resp. násobení) dvoch prirodzených čísel si preto všímame párnosť poslednej číslice. V tabuľke je prehľad súčtov a súčinov všetkých jednociferných prirodzených čísel.

\begin{table}[H]
  \centering
  \renewcommand{\arraystretch}{1.5}
  \newcolumntype{C}{>{\centering\arraybackslash}m{1em}}
  \begin{tabular}{C|CCCCC}
    \(+\) & 0 & 2 & 4 & 6 & 8 \\
    \hline
    0 & 0 & 2 & 4 & 6 & 8 \\
    2 & 2 & 4 & 6 & 8 & 10 \\
    4 & 4 & 6 & 8 & 10 & 12 \\
    6 & 6 & 8 & 10 & 12 & 14 \\
    8 & 8 & 10 & 12 & 14 & 16 \\
  \end{tabular}
  \hspace{1em}
  \begin{tabular}{C|CCCCC}
    \(+\) & 1 & 3 & 5 & 7 & 9 \\
    \hline
    1 & 2 & 4 & 6 & 8 & 10 \\
    3 & 4 & 6 & 8 & 10 & 12 \\
    5 & 6 & 8 & 10 & 12 & 14 \\
    7 & 8 & 10 & 12 & 14 & 16 \\
    9 & 10 & 12 & 14 & 16 & 18 \\
  \end{tabular}
  \hspace{1em}
  \begin{tabular}{C|CCCCC}
    \( + \) & 1 & 3 & 5 & 7 & 9 \\
    \hline
    0 & 1 & 3 & 5 & 7 & 9 \\
    2 & 3 & 5 & 7 & 9 & 11 \\
    4 & 5 & 7 & 9 & 11 & 13 \\
    6 & 7 & 9 & 11 & 13 & 15 \\
    8 & 9 & 11 & 13 & 15 & 17 \\
  \end{tabular}

  \vspace{1em}
  \begin{tabular}{C|CCCCC}
    \(\cdot\) & 0 & 2 & 4 & 6 & 8 \\
    \hline
    0 & 0 & 0 & 0 & 0 & 0 \\
    2 & 0 & 4 & 8 & 12 & 16 \\
    4 & 0 & 8 & 16 & 24 & 32 \\
    6 & 0 & 12 & 24 & 36 & 48 \\
    8 & 0 & 16 & 32 & 48 & 64 \\
  \end{tabular}
  \hspace{1em}
  \begin{tabular}{C|CCCCC}
    \(\cdot\) & 1 & 3 & 5 & 7 & 9 \\
    \hline
    1 & 1 & 3 & 5 & 7 & 9 \\
    3 & 3 & 9 & 15 & 21 & 27 \\
    5 & 5 & 15 & 25 & 35 & 45 \\
    7 & 7 & 21 & 35 & 49 & 63 \\
    9 & 9 & 27 & 45 & 63 & 81 \\
  \end{tabular}
  \hspace{1em}
  \begin{tabular}{C|CCCCC}
    \(\cdot\) & 1 & 3 & 5 & 7 & 9 \\
    \hline
    0 & 0 & 0 & 0 & 0 & 0 \\
    2 & 2 & 6 & 10 & 14 & 18 \\
    4 & 4 & 12 & 20 & 28 & 36 \\
    6 & 6 & 18 & 30 & 42 & 54 \\
    8 & 8 & 24 & 40 & 56 & 72 \\
  \end{tabular}
  \caption{Prehľad súčtov a súčinov párnych a nepárnych cifier.}
\end{table}

Na základe všetkých kombinácií sčítania a násobenia jednociferných prirodzených čísel vieme vyvodiť nasledovné vlastnosti párnych a nepárnych čísel:
\begin{itemize}
  \item Súčet dvoch párnych čísel je párne číslo. Ak sčítame ľubovoľné dve párne čísla, číslica na mieste jednotiek súčtu závisí od súčtu číslic rádu jednotiek dvoch sčítancov. Párne čísla \(0\), \(2\), \(4\), \(6\) a \(8\) vždy dávajú párny súčet.
  \item Súčet dvoch nepárnych čísel je párne číslo. Ak sčítame ľubovoľné dve nepárne čísla, číslica na mieste jednotiek súčtu závisí od súčtu číslic rádu jednotiek dvoch sčítancov. Nepárne čísla \(1\), \(3\), \(5\), \(7\) a \(9\) vždy dávajú párny súčet.
  \item Súčet párneho a nepárneho čísla je nepárne číslo. Ak sčítame ľubovoľné párne a nepárne číslo, číslica na mieste jednotiek súčtu závisí od súčtu číslic rádu jednotiek dvoch sčítancov. Súčet párnych čísel \(0\), \(2\), \(4\), \(6\) a \(8\) a nepárnych čísel \(1\), \(3\), \(5\), \(7\) a \(9\) vždy dáva nepárny súčet.
  \item Súčin dvoch párnych čísel je párne číslo. Ak vynásobíme ľubovoľné dve párne čísla, číslica na mieste jednotiek súčinu závisí od súčinu číslic rádu jednotiek dvoch činiteľov. Párne čísla \(0\), \(2\), \(4\), \(6\) a \(8\) vždy dávajú párny súčin.
  \item Súčin dvoch nepárnych čísel je nepárne číslo. Ak vynásobíme ľubovoľné dve nepárne čísla, číslica na mieste jednotiek súčinu závisí od súčinu číslic rádu jednotiek dvoch činiteľov. Nepárne čísla \(1\), \(3\), \(5\), \(7\) a \(9\) vždy dávajú nepárny súčin.
  \item Súčin párneho a nepárneho čísla je párne číslo. Ak vynásobíme ľubovoľné párne a nepárne číslo, číslica na mieste jednotiek súčinu závisí od súčinu číslic rádu jednotiek dvoch činiteľov. Súčin párnych čísel \(0\), \(2\), \(4\), \(6\) a \(8\) a nepárnych čísel \(1\), \(3\), \(5\), \(7\) a \(9\) vždy dáva párny súčin.
\end{itemize}
Zhrnúť tieto vlastnosti môžeme nasledove. Súčet dvoch čísel rovnakej párnosti je párny. Súčet dvoch čísel rôznej párnosti je nepárny. Ak je aspoň jeden činiteľ párny, súčin je párny. Súčin je nepárny iba vtedy, keď sú oba činitele nepárne. V tabuľke sú znázornené výsledky sčítania a násobenia párnych a nepárnych čísel.



\begin{table}[H]
  \centering
  \renewcommand{\arraystretch}{1.5}
  \begin{tabular}{c|cc}
    \(+\) & Párne & Nepárne\\
    \hline
    Párne & Párne & Nepárne \\
    Nepárne & Nepárne & Párne \\
  \end{tabular}
  \hspace{1em}
  \begin{tabular}{c|cc}
    \(\cdot\) & Párne & Nepárne\\
    \hline
    Párne & Párne & Párne \\
    Nepárne & Párne & Nepárne \\
  \end{tabular}
  \caption{Prehľad výsledkov sčítania a násobenia párnych a nepárnych čísel.}
\end{table}

\section{Prvočísla}

\section{Dokonalé čísla}

Dokonalé číslo je celé číslo, ktoré je súčtom všetkých svojich deliteľov okrem seba samého.
Prvých desať dokonalých čísel:

6
28
496
8128
33 550 336
8 589 869 056
137 438 691 328
2 305 843 008 139 952 128
2 658 455 991 569 831 744 654 692 615 953 842 176
191 561 942 608 236 107 294 793 378 084 303 638 130 997 321 548 169 216


\end{document}